% ─────────────────────────────────────────────────────────────────────────────
% Paper 2 — Symbolic Substrates (VSL inter-LLM protocol) — SEED
% Extracted verbatim from attractor_convergence/main.tex (2026-06-14, step 10 of
% the narrative revision). This is Mode-2 material (the semantic substrate) and is
% off-thesis for Paper 1 (the attractor); it is the natural empirical anchor of
% Paper 2 (see memory project_paper2_seed, EXP-PHI-01).
%
% Paper 1 retains ONE sentence pointing here (sec:vsl) + a 1-line pointer in the
% Maxwell/Observer experiment ("VSL as the Demon's apparatus").
%
% The structural≈100% vs semantic≈40% drift finding + the domain_context anchor
% are the core of Paper 2: a viability measurement of Mode 2 (does meaning survive
% the substrate crossing?).
% ─────────────────────────────────────────────────────────────────────────────

\section{VSL: Semantic Compression as Inter-Layer Protocol}
\label{sec:vsl}

\subsection*{Question}
Can a compact formal language preserve operational semantics across system
boundaries with significantly fewer tokens than natural language?

\subsection{Motivation}

The experiments described above required communicating complex protocol
specifications -- thermal regimes, stability thresholds, governance axioms --
between system layers and across agent boundaries.
Natural language instructions are verbose, ambiguous, and sensitive to
paraphrase; they do not compress cleanly.

\emph{Viable System Language} (VSL) is a machine-readable notation derived
from Beer's VSM formalism.
A VSL instruction expresses system hierarchy, axioms, state transitions,
and metrics in a grammar that can be parsed without prose disambiguation.
We measured whether this compactness is real and whether it survives
cross-agent transmission.

\subsection{Protocol}

Five protocol specifications of increasing complexity were written twice:
once as natural-language prompts (\texttt{prompt\_A}), once as VSL
instructions (\texttt{instruccion\_B}).
The B versions were further split into two variants:
\texttt{B\_full} (standalone -- recipient has no prior schema)
and \texttt{B\_schema} (machine-to-machine -- recipient has VSF schema
pre-loaded).
Token counts used \texttt{tiktoken} (\texttt{cl100k\_base}).

\subsection{Results}

\begin{table}[h]
\centering
\caption{VSL compression results (tiktoken cl100k\_base).
  $r_\text{full}$ = standalone compression ratio;
  $r_\text{schema}$ = machine-to-machine compression ratio.}
\label{tab:vsl}
\begin{tabular}{lrrrrr}
\hline
Case & $T_A$ & $T_{B,\text{full}}$ & $T_{B,\text{schema}}$
     & $r_\text{full}$ & $r_\text{schema}$ \\
\hline
\texttt{pvsnp}               & 328 & 161 &  --  & 0.51 &  --  \\
\texttt{thermal\_senescence} & 638 & 346 &  --  & 0.46 &  --  \\
\texttt{spec\_boot+Lyapunov} & 717 & 432 & 148  & 0.40 & 0.79 \\
\texttt{proxmox+S4}          & 698 & 411 & 198  & 0.41 & 0.72 \\
\texttt{PSI/1.0}             & 641 & 397 & 147  & 0.38 & 0.77 \\
\hline
\end{tabular}
\end{table}

Standalone VSL achieves $38$--$51\%$ compression.
With a shared schema (machine-to-machine mode), compression rises to
$72$--$79\%$ and grows with protocol complexity, consistent with the
original $83\%$ observation on a single case.

\subsection{Semantic drift test}

The five \texttt{B\_schema} instructions were submitted to an independent
language model (Gemini) with the VSF schema pre-loaded.
Structural fidelity -- correct identification of S-layer hierarchy,
$\varphi\!\uparrow/\!\downarrow$ channels, PSI thresholds, covenant
verification step, and metric names -- was $\approx 100\%$.

Semantic fidelity was $\approx 40\%$: the model systematically substituted
its own operational context for VSF domain referents
(e.g.\ $\varphi\!\uparrow$ checkpoints mapped to cloud-sync operations;
\texttt{vm-nucleus} mapped to an internal knowledge module).
This \emph{context substitution} is structurally silent -- the parse tree
is correct but the referents are wrong.

The fix is a lightweight domain anchor:

\begin{verbatim}
domain_context: {target=vsf-nucleus, env=Rust/LXC, layer=S1}
\end{verbatim}

\noindent added to the \texttt{B\_schema} header.
This adds $\approx 8$ tokens and eliminates context substitution in
subsequent tests.

\textbf{Cross-model replication.}
As an independent replication, the VSL v1.0 specification was provided to
Qwen3.6-Plus (a distinct architecture and training distribution from Gemini
and Claude).
Without additional instructions, the model parsed the grammar and produced
a syntactically valid \texttt{spec.vsm} document following all VSL
conventions.
This constitutes a third independent model confirming structural fidelity,
consistent with the hypothesis that VSL grammar is learnable from the
specification alone.

\subsection{Hypotheses result}

\begin{description}
  \item[H1 -- confirmed.]  Schema mode: $72$--$79\% \geq 70\%$ threshold.
  \item[H2 -- partial.]    Structure preserved; domain grounding requires
        the anchor block.
  \item[H3 -- split.]      Structural drift $\approx 0\%$;
        semantic drift $\approx 60\%$ without anchor, $\approx 5\%$ with
        anchor.  The original single metric conflated two independent
        properties.
\end{description}

\subsection{Implication}

VSL compression enables governance signals to cross organisational
boundaries -- between kernels running on different hardware, maintained
by different teams, implemented in different languages -- at a token cost
that scales with protocol complexity rather than prose verbosity.
The structural/semantic drift distinction introduced here is a prerequisite
for any formal correctness guarantee of inter-agent communication.
